\documentclass[11pt,a4paper]{article}
\usepackage[utf8]{inputenc}
\usepackage[T1]{fontenc}
\usepackage[spanish,es-tabla]{babel}
\usepackage{lmodern,microtype,geometry,graphicx,amsmath,amssymb,amsthm,booktabs,array,xcolor,float,hyperref}
\geometry{margin=2.6cm}
\hypersetup{colorlinks=true,linkcolor=blue!60!black,urlcolor=blue!60!black,citecolor=teal!50!black}
\newtheorem{proposicion}{Proposición}
\newcommand{\um}{\,\mu\mathrm{g}\,\mathrm{m}^{-3}}
\newcommand{\real}{\mathbb R}
\newcommand{\RRMSE}{\mathrm{RMSE}}
\graphicspath{{../results/figures/}}
\title{\Huge\textbf{NEXO}\\[.4em]\Large Reconstrucción espacial, análisis vectorial\\ e inferencia estadística de campos ambientales urbanos\\[.8em]\normalsize Prototipo matemático reproducible y verificación mediante datos sintéticos controlados}
\author{Laboratorio interdisciplinario de reconstrucción matemática\\ y modelación ambiental urbana}
\date{18 de octubre de 2026 \quad Versión 0.1}
\begin{document}
\maketitle\vspace{1cm}
\begin{center}\fcolorbox{orange!65!black}{orange!9}{\parbox{.85\textwidth}{\centering\textbf{Advertencia sobre los datos.} Los números y mapas de resultados de este documento se obtuvieron exclusivamente de campos y estaciones \textbf{sintéticos}. No se han descargado mediciones reales en el entorno de ejecución; no se presentan conclusiones empíricas sobre Monterrey.}}\end{center}
\vspace{1cm}
\begin{abstract}
Se desarrolla una primera implementación reproducible de \emph{NEXO}, un laboratorio de reconstrucción de campos escalares a partir de observaciones espaciales escasas. Se contrastan ponderación por distancia inversa y bases radiales gaussianas regularizadas, empleando descomposición en valores singulares, validación dejando fuera una estación, diferenciación espacial y un esquema conservativo para la ecuación de advección--difusión. Se identifican fuentes de observaciones ambientales y productos meteorológicos para el área metropolitana de Monterrey, pero el acceso a las mediciones no pudo completarse desde el entorno disponible. Por consiguiente, todas las cifras experimentales corresponden a casos sintéticos con función verdadera conocida, semilla declarada y generación automática de figuras. Se discuten error de generalización, efecto del condicionamiento y pruebas de estabilidad, convergencia y conservación. Los resultados respaldan la corrección del prototipo bajo hipótesis controladas; no constituyen una validación del transporte atmosférico de Monterrey.
\end{abstract}
\noindent\textbf{Palabras clave:} problema inverso; interpolación espacial; regularización de Tikhonov; SVD; cálculo vectorial; advección--difusión; validación espacial; calidad del aire.\\[1em]
\tableofcontents\clearpage
\input{sections/01_contexto}
\input{sections/02_matematicas}
\input{sections/03_experimentos}
\input{sections/04_discusion}
\input{sections/05_reproducibilidad}
\bibliographystyle{plain}
\bibliography{references}
\appendix
\input{appendix/algoritmos}
\end{document}
